\chapter*{Annexes}
\addcontentsline{toc}{chapter}{Annexes}
\markboth{Annexes}{Annexes}

% ============================================================
\section*{Annexe A --- Code Python du projet}
\label{ann:code}
% ============================================================

\begin{lstlisting}[caption={Correction des dates flottantes (data cleaning)},
                   label={lst:date_fix}]
def fix_date(val):
    """Corrige les dates au format flottant (20100312.0 -> 2010-03-12)."""
    if pd.isna(val):
        return pd.NaT
    s = str(val).split('.')[0]  # supprime la partie decimale
    try:
        return pd.to_datetime(s, format='%Y%m%d')
    except ValueError:
        return pd.NaT

df['BENE_BIRTH_DT'] = df['BENE_BIRTH_DT'].apply(fix_date)
df['BENE_DEATH_DT'] = df['BENE_DEATH_DT'].apply(fix_date)
\end{lstlisting}

\begin{lstlisting}[caption={Winsorisation des colonnes de coûts},
                   label={lst:winsor}]
cost_cols = [
    'MEDREIMB_IP', 'BENRES_IP', 'PPPYMT_IP',
    'MEDREIMB_OP', 'BENRES_OP', 'PPPYMT_OP',
    'MEDREIMB_CAR', 'BENRES_CAR', 'PPPYMT_CAR'
]
for col in cost_cols:
    p99 = df[col].quantile(0.99)
    df[col] = df[col].clip(upper=p99)
\end{lstlisting}

\begin{lstlisting}[caption={Calcul de l'index de Charlson},
                   label={lst:charlson}]
def compute_charlson(row):
    """Index de Charlson pondere par comorbidite."""
    score = 0
    score += 2 * row['SP_CHRNKIDN']   # maladie renale (poids 2)
    score += 2 * row['SP_CNCR']        # cancer (poids 2)
    score += 2 * row['SP_STRKETIA']    # AVC/AIT (poids 2)
    for col in ['SP_CHF', 'SP_DIABETES', 'SP_COPD',
                'SP_ALZHDMTA', 'SP_DEPRESSN',
                'SP_ISCHMCHT', 'SP_RA_OA']:
        score += row[col]              # poids 1
    return score

df['CHARLSON_INDEX'] = df.apply(compute_charlson, axis=1)
\end{lstlisting}

\begin{lstlisting}[caption={Optimisation XGBoost avec Optuna et logging MLflow},
                   label={lst:optuna}]
import optuna, mlflow
from xgboost import XGBClassifier
from sklearn.metrics import roc_auc_score

def objective(trial):
    params = {
        'n_estimators'     : trial.suggest_int('n_estimators', 100, 600),
        'max_depth'        : trial.suggest_int('max_depth', 3, 10),
        'learning_rate'    : trial.suggest_float('learning_rate', 0.01, 0.2),
        'subsample'        : trial.suggest_float('subsample', 0.6, 1.0),
        'colsample_bytree' : trial.suggest_float('colsample_bytree', 0.5, 1.0),
        'scale_pos_weight' : 10.94,
        'tree_method'      : 'hist',
        'eval_metric'      : 'auc',
        'use_label_encoder': False,
    }
    model = XGBClassifier(**params)
    model.fit(X_train, y_train,
              eval_set=[(X_val, y_val)],
              verbose=False)
    return roc_auc_score(y_val, model.predict_proba(X_val)[:,1])

study = optuna.create_study(direction='maximize')
study.optimize(objective, n_trials=50)

# Logging MLflow du meilleur run
with mlflow.start_run(run_name="xgboost_optuna_best"):
    mlflow.log_params(study.best_params)
    mlflow.log_metric("auc_val",   0.9245)
    mlflow.log_metric("auc_test",  0.8973)
    mlflow.log_metric("recall",    0.9150)
    mlflow.log_metric("f1_score",  0.4210)
    mlflow.log_metric("threshold", 0.87)
    mlflow.xgboost.log_model(best_model, "model",
                             registered_model_name="HospitalisationRisk")
\end{lstlisting}


\begin{lstlisting}[caption={Endpoint FastAPI /predict avec schéma Pydantic},
                   label={lst:predict}]
from fastapi import FastAPI
from pydantic import BaseModel
import mlflow.pyfunc

app = FastAPI(title="Hospitalization Risk API", version="3.0.0")
model = mlflow.pyfunc.load_model("models:/HospitalisationRisk/Production")

class PatientFeatures(BaseModel):
    AGE: int; SEXE_ENC: int; RACE_ENC: int
    BENE_ESRD_IND: int; GROUPE_AGE_ENC: int
    SP_ALZHDMTA: int; SP_CHF: int; SP_CHRNKIDN: int
    SP_CNCR: int; SP_COPD: int; SP_DEPRESSN: int
    SP_DIABETES: int; SP_ISCHMCHT: int; SP_OSTEOPRS: int
    SP_RA_OA: int; SP_STRKETIA: int
    NB_COMORBIDITES: float; CHARLSON_INDEX: float
    COUT_TOTAL: float; IS_NEW_PATIENT: int
    NB_HOSP_PASSEES: int; NB_OP_3M: int; NB_OP_6M: int
    NB_OP_12M: int; NB_CAR_6M: int; NB_PRESCRIPTIONS: int
    NB_MOLECULES_UNIQUES: int; POLYPHARMACIE: int
    SP_STATE_CODE: int

@app.post("/predict")
async def predict(patient: PatientFeatures):
    import pandas as pd
    df = pd.DataFrame([patient.dict()])
    proba = float(model.predict(df)[0])
    if proba < 0.40:
        niveau = "FAIBLE"
    elif proba < 0.70:
        niveau = "MODERE"
    else:
        niveau = "ELEVE"
    return {"score_risque": round(proba, 4),
            "niveau_risque": niveau,
            "seuil_decision": 0.87}
\end{lstlisting}

\begin{lstlisting}[caption={Flow Prefect de réentraînement automatique},
                   label={lst:prefect}]
from prefect import flow, task
from datetime import timedelta

@task(retries=3, retry_delay_seconds=60)
def validate_data():
    return run_ge_checkpoint("hospitalization_suite")

@task
def retrain_model(data_version: str):
    return run_training_pipeline(data_version)

@task
def evaluate_and_register(model_path: str):
    auc = evaluate_model(model_path)
    if auc > 0.85:
        mlflow.register_model(model_path, "HospitalisationRisk")
    return auc

@flow(name="hospitalization-retraining",
      schedule=timedelta(weeks=4))
def retraining_flow():
    data_ok = validate_data()
    if data_ok:
        model = retrain_model("latest")
        auc = evaluate_and_register(model)
        if auc < 0.85:
            rollback_to_previous_version()
\end{lstlisting}

\begin{lstlisting}[caption={Expectations Great Expectations},
                   label={lst:ge}]
import great_expectations as ge

df_ge = ge.from_pandas(df)

# Contrainte 1 : nombre de lignes
df_ge.expect_table_row_count_to_be_between(
    min_value=85_000, max_value=95_000)

# Contrainte 2 : colonnes requises
df_ge.expect_column_to_exist('CHARLSON_INDEX')
df_ge.expect_column_to_exist('NB_COMORBIDITES')

# Contrainte 3 : taux d'hospitalisation coherent
df_ge.expect_column_mean_to_be_between(
    'HOSP_ANNEE_SUIV', min_value=0.06, max_value=0.12)
\end{lstlisting}

% ============================================================
\section*{Annexe B --- Commandes essentielles}
\label{ann:commandes}
% ============================================================

\begin{lstlisting}[language=bash, caption={Commandes principales du projet},
                   label={lst:commands}]
# ---- Environnement ----
python -m venv .venv && source .venv/bin/activate
pip install -r requirements.txt

# ---- DVC ----
dvc pull                         # telecharger les donnees
dvc add data/cleaned/
git add data/cleaned.dvc && git commit -m "feat: cleaned data v1"
dvc repro                        # rejouer le pipeline complet

# ---- Great Expectations ----
great_expectations checkpoint run hospitalization_suite

# ---- Entrainement ----
python src/train.py --model xgboost --n_trials 50

# ---- MLflow UI ----
mlflow ui --port 5000

# ---- Docker ----
docker compose up -d             # demarrer tous les services
docker compose down              # arreter les services
docker compose logs -f api       # logs de l'API

# ---- Tests ----
pytest tests/ -v --cov=src --cov-report=html

# ---- Monitoring ----
python scripts/compute_psi.py --reference 2008 --current 2010

# ---- API (test manuel) ----
curl -X POST http://localhost:8000/predict \
  -H "Content-Type: application/json" \
  -d '{"AGE":72,"NB_COMORBIDITES":4,"CHARLSON_INDEX":5.0,
       "NB_HOSP_PASSEES":2,"NB_OP_6M":8,"COUT_TOTAL":45000,
       "SEXE_ENC":0,"RACE_ENC":1,"BENE_ESRD_IND":0,
       "GROUPE_AGE_ENC":3,"SP_CHF":1,"SP_DIABETES":1,
       "SP_ALZHDMTA":0,"SP_RA_OA":0,"SP_ISCHMCHT":1,
       "SP_OSTEOPRS":0,"SP_COPD":0,"SP_DEPRESSN":0,
       "SP_STRKETIA":0,"SP_CNCR":0,"SP_CHRNKIDN":0,
       "NB_OP_3M":2,"NB_OP_12M":8,"NB_CAR_6M":3,
       "NB_PRESCRIPTIONS":7,"NB_MOLECULES_UNIQUES":6,
       "POLYPHARMACIE":1,"IS_NEW_PATIENT":0,"SP_STATE_CODE":10}'
\end{lstlisting}

% ============================================================
\section*{Annexe C --- requirements.txt}
\label{ann:requirements}
% ============================================================

\begin{lstlisting}[language=bash, caption={Dépendances Python du projet},
                   label={lst:requirements}]
# Data
pandas==2.1.4
numpy==1.26.3
pyarrow==14.0.2

# Machine Learning
scikit-learn==1.4.0
xgboost==2.0.3
lightgbm==4.2.0
optuna==3.5.0
shap==0.44.0

# MLOps
mlflow==2.10.0
dvc==3.38.1
great-expectations==0.18.8
prefect==2.14.21

# API
fastapi==0.109.0
uvicorn[standard]==0.27.0
pydantic==2.5.3

# Monitoring
evidently==0.4.17
grafana-client==3.7.0

# Visualization
matplotlib==3.8.2
seaborn==0.13.1

# Tests
pytest==7.4.4
pytest-cov==4.1.0
httpx==0.26.0
\end{lstlisting}

% ============================================================
\section*{Annexe D --- Métriques de performance complètes}
\label{ann:metriques}
% ============================================================

\begin{table}[htbp]
\centering
\caption{Métriques complètes sur les trois jeux de données}
\label{tab:metriques_completes}
\begin{tabular}{@{}lllllll@{}}
\toprule
\textbf{Modèle} & \textbf{Jeu} & \textbf{AUC} & \textbf{Recall} &
\textbf{Précision} & \textbf{F1} & \textbf{$\tau^*$} \\
\midrule
\multirow{3}{*}{Logistique}
  & Train 2008 & 0.9214 & 0.9512 & 0.1887 & 0.3164 & \multirow{3}{*}{0.961} \\
  & Val. 2009  & 0.9082 & 0.9339 & 0.2401 & 0.3825 & \\
  & Test 2010  & 0.8847 & 0.9121 & 0.2109 & 0.3431 & \\
\midrule
\multirow{3}{*}{XGBoost}
  & Train 2008 & 0.9876 & 0.9810 & 0.3241 & 0.4880 & \multirow{3}{*}{0.870} \\
  & Val. 2009  & 0.9245 & 0.9150 & 0.2913 & 0.4210 & \\
  & Test 2010  & 0.8973 & 0.8920 & 0.2540 & 0.3942 & \\
\midrule
\multirow{3}{*}{LightGBM}
  & Train 2008 & 0.9843 & 0.9790 & 0.3108 & 0.4712 & \multirow{3}{*}{0.830} \\
  & Val. 2009  & 0.9210 & 0.9200 & 0.2871 & 0.4150 & \\
  & Test 2010  & 0.8934 & 0.8970 & 0.2481 & 0.3885 & \\
\bottomrule
\end{tabular}
\end{table}

